\section{The universal defect and its coproduct}
\label{cp:sec-relative-rigidity}

Fix a physical point on one of the two charts supplied by
Proposition~\ref{cp:prop-physical-point-on-chart}.  Shrink its connected
principal open \(U\), without removing that point, so that all matrices,
cross-ratios, and finite root covers mentioned below are regular.  This
does not change the function field
\(\Q(t_1,\ldots,t_6)\) of the chart.

\subsection{The matrix geometry behind the motivic input}

Let \(G\) and \(G^R\) be the universal angle Gram matrices with their common
determinant root \(w\).  We work with the projectively dual coordinate
quadrics
\[
 Q_G:\ z^tGz=0,\qquad Q_{G^R}:\ z^tG^Rz=0.
\]
The following calculation explains the shape of the three-weight package.
It does not construct that package.

\begin{lemma}[Coordinate-stratum calculation]
\label{cp:lem-split-strata}
At the level of quadratic forms, \(Q_G\) and \(Q_{G^R}\) are hyperbolic on
\(U\); every coordinate rank-three restriction has an isotropic line.  Five
coordinate binary restrictions have their two isotropic points displayed by
the phases \(a,b,c,d,q\).  The sixth has isotropic points over the common
quadratic cover
\[
 U_p=U[p]/(p^2+2xp+1),
\]
whose deck transformation is \(p\mapsto p^{-1}\).
\end{lemma}

\begin{proof}
For a unit \(u\), put \(g=-(u+u^{-1})/2\).  Then \((1,u)\) is isotropic
for
\[
 \begin{pmatrix}1&g\\g&1\end{pmatrix},
\qquad 1+2gu+u^2=0.
\]
In particular, the \(q\)-block is a hyperbolic plane.  Its binary Schur
complement \(C\) in \(G\) satisfies
\[
 \det C=\frac{\det G}{1-y^2}
       =-\left(\frac ws\right)^2,
\qquad s=\frac{q-q^{-1}}2.
\]
Put \(t=w/s\), \(J_0=\left(\begin{smallmatrix}0&1\\-1&0\end{smallmatrix}\right)\),
and \(J=t^{-1}J_0C\).  The identities
\[
 (J_0C)^2=-\det(C)I,\qquad J^2=I,\qquad J^tC+CJ=0
\]
show that the two eigenlines of \(J\) are isotropic and paired perfectly by
\(C\).  Thus the complementary plane, and hence the rank-four form, is
hyperbolic.  The same \(q,s,w\) are fixed by Regge, so the calculation also
applies to \(G^R\).

Every coordinate rank-three restriction contains one of the displayed
binary blocks and hence an isotropic vector.  The five phase variables give
the two roots of five binary restrictions.  The remaining restriction has
equation \(p+p^{-1}=-2x\), equivalently \(p^2+2xp+1=0\), and
Proposition~\ref{cp:prop-fixed-p-cover} gives its deck transformation.
\end{proof}

Passing from this calculation to relative motives of all strata, including
regularity, functorial localization, frames, and the reduced sign-Artin
summand, is part of Inputs \textup{(E2)--(E4)}.  We do not relabel that
passage as a proof by localization.

\subsection{The relative coefficient supplied by the contract}

Apply Input~\textup{(E2)} to the two oriented families.  In the notation of
Definition~\ref{cp:def-coefficient-contract}, set
\[
 c_G=c_U(G),\qquad c_R=c_U(G^R),\qquad
 \delta_U^\vee=c_G-c_R\in A_{02}(U).
\]
The superscript \(\vee\) records projective duality
\(Q_{G^{-1}}\leftrightarrow Q_G\); it is not the internal dual used in the
intended face motive.

Input~\textup{(E3)} supplies the following statement.  It is displayed as
an external proposition so that its logical status cannot be mistaken.

\begin{proposition}[External three-weight package]
\label{cp:prop-weight-calculation}
The two coefficients have endpoint types \((0,2)\), their middle terms are
five Tate channels and one possible sign-Artin channel, and
\[
 \operatorname{prim}_U:
 \ker\overline\Delta_U\xrightarrow{\sim}H^1(U,\Q(2))
\]
is an isomorphism.  Moreover \(H^1(U,\Q(2))=0\).
\end{proposition}

This proposition is Input~\textup{(E3)}, not a theorem proved from the
coordinate-stratum calculation.  Its proposed motivic graded pieces are
\[
 \Q_U(0),\qquad \Q_U(1)^{\oplus5}\oplus\chi(1),\qquad\Q_U(2),
\]
and its proposed vanishing route is Equation
\eqref{cp:eq-weight-two-rigidity-route}.

\subsection{Length--angle channels}

Input~\textup{(E4)} identifies the reduced coproduct of a dual oriented
quadric coefficient with
\begin{equation}
 \overline\Delta_U c_G=\sum_e L_e\otimes A_e,
 \label{cp:eq-angle-length-coproduct}
\end{equation}
one term for each physical edge.  The angle factor comes from a coordinate
binary section and the length factor from its complementary binary Schur
form.  This ordering is important: it is \(L_e\otimes A_e\).

Here is the elementary cross-ratio normalization.  For
\[
 r(z_1,z_2,z_3,z_4)
 =\frac{(z_1-z_3)(z_2-z_4)}
        {(z_1-z_4)(z_2-z_3)}
\]
and two real rays of angle \(\phi\), use slopes \(0,\tan\phi\) and order the
isotropic directions as \(-\mathrm i,\mathrm i\).  Direct substitution gives
\[
 r(-\mathrm i,\mathrm i,0,\tan\phi)
 =\frac{1+\mathrm i\tan\phi}{1-\mathrm i\tan\phi}
 =e^{2\mathrm i\phi}.
\]
Thus, with compatible orientations,
\begin{equation}
 A_e=e^{2\mathrm i\vartheta_e},
 \qquad L_e=e^{2\mathrm i\ell_e}.
 \label{cp:eq-cross-ratio-phases}
\end{equation}
Swapping the two isotropic directions inverts the cross-ratio.  Input
\textup{(E4)} asserts that its relative formula uses exactly these
orientations and that the same convention survives specialization.

\subsection{The internal tensor calculation}

Let
\[
 R=\frac12J_4-I_4,\qquad H=2R=J_4-2I_4.
\]
Section~\ref{cp:sec-phase-charts} proves \(R^tR=I_4\), hence
\begin{equation}
 H^tH=4I_4.
 \label{cp:eq-H-orthogonal}
\end{equation}
For the four moving edges, choose half-phases
\[
 A_i=B_i^2,\qquad L_i=Y_i^2.
\]
The universalization clause of Input~\textup{(E4)} says that, on a common
finite splitting cover, the real Regge relations of Input~\textup{(E1)}
extend to the algebraic identities
\begin{equation}
 A_i^R=\prod_jB_j^{H_{ij}},
 \qquad
 L_i^R=\prod_jY_j^{H_{ij}}.
 \label{cp:eq-monomial-angle-length}
\end{equation}
This clause includes the Zariski-density argument and the algebraic recovery
of length phases from cofactors; neither is silently inferred from the
formal matrix identity.

\begin{lemma}[Regge tensor cancellation]
\label{cp:lem-regge-tensor-cancellation}
Let \(B_1,\ldots,B_4\) and \(Y_1,\ldots,Y_4\) lie in abelian multiplicative
groups, tensor those groups over \(\mathbf Z\), and then tensor with
\(\Q\).  If \(A_i=B_i^2\), \(L_i=Y_i^2\), and the transformed terms obey
\eqref{cp:eq-monomial-angle-length}, then
\[
 \sum_i L_i^R\otimes A_i^R=\sum_iL_i\otimes A_i.
\]
\end{lemma}

\begin{proof}
Write a multiplicative group additively inside its tensor product.  Then
\([u^m]=m[u]\), so Equation~\eqref{cp:eq-monomial-angle-length} and
\(H^tH=4I_4\) give
\[
 \begin{aligned}
 \sum_i L_i^R\otimes A_i^R
 &=\sum_{i,j,k}H_{ik}H_{ij}\,[Y_k]\otimes[B_j]\\
 &=4\sum_j[Y_j]\otimes[B_j]\\
 &=\sum_j[Y_j^2]\otimes[B_j^2]
  =\sum_jL_j\otimes A_j.
 \end{aligned}
\]
Changing a square root changes its class by two-torsion, which disappears
after tensoring with \(\Q\).
\end{proof}

This is the entire Regge-specific coproduct computation.  Lean proves both
this half-phase form using \(H\) and the equivalent full-channel form using
the normalized matrix \(R=H/2\).  The latter is the identity connected
directly to the six-channel coproduct contract.

\subsection{Primitive and universal vanishing}

\begin{lemma}[Vanishing of the middle coproduct]
\label{cp:lem-coproduct-zero}
Under Inputs \textup{(E1), (E2), and (E4)},
\[
 \overline\Delta_U(\delta_U^\vee)=0.
\]
At every complex point \(s:\Spec\C\to U\), the specialized dual
quadric-scissors defect is primitive as well.
\end{lemma}

\begin{proof}
The two channels belonging to the selected opposite edge pair agree
term-by-term because Regge fixes the corresponding lengths and angles.  On
the common splitting cover, the remaining four channels agree by
Lemma~\ref{cp:lem-regge-tensor-cancellation}.  The descent clause of Input
\textup{(E4)} says that pullback is injective on the five Tate channels and,
using normalized twisted trace, on the possible sign-Artin channel.  Hence
the equality descends to \(A_{11}(U)\).  The coproduct-specialization square
in the same input carries this zero to Goncharov's field-level coproduct.
Thus the specialized defect lies in its kernel; this is exactly the proof
term which licenses applying \(c_G\), rather than an independently assumed
primitive label.
\end{proof}

\begin{theorem}[Vanishing of the universal framed defect]
\label{cp:thm-relative-defect-zero}
Under Inputs \textup{(E1)--(E4)},
\[
 \delta_U^\vee=0\qquad\text{in }A_{02}(U).
\]
\end{theorem}

\begin{proof}
Lemma~\ref{cp:lem-coproduct-zero} makes the defect primitive.
Proposition~\ref{cp:prop-weight-calculation} supplies the weight-two
coefficient contract.  Lemma~\ref{cp:lem-three-level-framed} therefore
gives \(\delta_U^\vee=0\).
\end{proof}

The proof of this theorem is now transparently divided.  Orthogonality and
the tensor cancellation are internal and machine checked.  Construction of
the relative coefficients, the universalization and descent of their
coproduct, and primitive \(=\) motivic extension are the named external
clauses.
